\section{Evaluation}
\label{sec:eval}

We evaluate \Name across five stages: static scarcity~(\S\ref{sec:eq1}), adversarial workload structure~(\S\ref{sec:eq2}), dynamic demand reversal~(\S\ref{sec:eq3}), trace-driven validation~(\S\ref{sec:trace}), and hardware validation~(\S\ref{sec:eq5}). Each stage uses a frozen scenario family (Table~\ref{tab:setup}); P-attn is the fraction of premium jobs with $\mathrm{SAS}_i \ge 0.98$. Static experiments run 600\,s; dynamic experiments evaluate three pre-registered windows; all numbers are mean $\pm$ std over ten seeds (five for the dynamic tier-swap and sensitivity experiments) unless noted otherwise. Additional robustness results are in Appendix~\ref{sec:app-robustness}.

\subsection{Experimental Setup}
\label{sec:setup}

\paragraph{Simulator.}
Our custom simulator models per-iteration training structure explicitly: each job alternates between computation (no network demand) and collective communication (duration depends on allocated bandwidth), so phase staggering~(\S\ref{sec:periodic-origin}) emerges from the model. SLO semantics follow \S\ref{sec:model}. We validate simulator fidelity against the physical testbed by replicating static scenarios bit-for-bit; the simulator reproduces the testbed's measured allocations and scheduler ranking (cross-validation in \S\ref{sec:impl-fidelity}).

\paragraph{Workloads, metrics, and baselines.}
Table~\ref{tab:setup} summarizes the four frozen scenario families and the primary metric. \textbf{E1} sweeps spine capacity across the feasibility boundary; \textbf{E2$'$} inverts the tier-to-model mapping; \textbf{E2-pro} is the positive control; \textbf{E3}/\textbf{E3$'$} execute a mid-run tier swap at $t{=}300$\,s. We compare against Fair, SRPT (shortest remaining processing time), CRUX~\cite{Crux}, CASSINI~\cite{Cassini}, DF (deficit-feedback, \S\ref{sec:why-closed-form}), and LL-S (LongLiu with a static $T$-target, \S\ref{sec:longliu}). P-attn is the fraction of premium jobs with $\mathrm{SAS}_i \ge 0.98$.

\begin{table}[h]
\centering
\caption{Frozen scenario families and primary metric.}
\label{tab:setup}
\footnotesize
\setlength{\tabcolsep}{3pt}
\begin{tabular}{@{}lcc@{}}
\toprule
\textbf{Scenario} & \textbf{Jobs} & \textbf{Spine (Gbps)} \\
\midrule
E1 (ladder) & 14 (8P/6S) & 400--1200 \\
E2$'$ (adversarial) & 14 (10P/4S) & 400--800 \\
E2-pro (favorable) & 14 (10P/4S) & 630, 800 \\
E3 (control) & 14 (8P/6S) & 800 \\
E3$'$ (kill) & 14 (8P/6S) & 630 \\
\bottomrule
\end{tabular}
\end{table}

\subsection{Static Scarcity}
\label{sec:eq1}

We first examine whether \Name differentiates from baselines under scarcity while remaining cost-free when bandwidth is abundant. \Name stratifies the field where bandwidth is scarce (Figure~\ref{fig:e1-ladder}). At 400~Gbps, \Name delivers $76.3\%$ P-attn vs.\ $73.8\%$ (best baseline, LL-S), a $2.5$~pp edge; at 500~Gbps the gap widens to $90.0\%$ vs.\ $80.0\%$ (SRPT, LL-S), a $10$~pp lead. At 630~Gbps, \Name reaches $96.3\%$ while the best baseline (Fair) stalls at $91.3\%$, and at 800~Gbps \Name ties DF at $97.5\%$. Above 800~Gbps all policies converge. On a separate 24-job anchor workload at 400~Gbps, \Name attains the highest SLO rate ($31.3\%$) among all seven policies (Table~\ref{tab:anchor} in Appendix~\ref{sec:app-supp}).

\begin{figure}[t]
\centering
\includegraphics[width=\columnwidth]{figure/fig2_e1_ladder.pdf}
\caption{Static scarcity: P-attn vs.\ spine bandwidth under E1 ($n{=}10$). \Name leads the next-best baseline by $10$~pp at 500~Gbps and converges with all baselines above 800~Gbps.}
\label{fig:e1-ladder}
\end{figure}

\paragraph{Generality across collectives.}
\Name schedules at the bandwidth-allocation layer: each job's communication is modeled only by its per-iteration volume $D_i$ and allocated bandwidth $a_i$ (comm.\ time $D_i/a_i$), an abstraction that holds for any bandwidth-limited collective---AllReduce (DDP), AllGather/ReduceScatter (ZeRO-3), and All-to-All (MoE expert routing) alike. To validate that the guarantees survive realistic mixed workloads, E17 replays E1's 14 jobs (8P/6S) with heterogeneous collectives: five AllReduce, four AllGather, three All-to-All, and two ReduceScatter jobs, with per-job communication volumes scaled to each collective ($0.5\times$ for the single-phase primitives). \Name remains the top policy at every scarce bandwidth: $67.5\%$ at 400~Gbps (vs.\ $61.3\%$ for Fair), $71.2\%$ at 500~Gbps, and $73.8\%$ at 630~Gbps, and it significantly beats DF ($p{=}0.006$), CRUX ($p{=}0.003$), CASSINI ($p{=}2{\times}10^{-4}$), and LL-S ($p{=}0.040$) at 400~Gbps (Figure~\ref{fig:e17-mixed}). At 800~Gbps all policies converge ($70$--$77.5\%$; no pairwise difference is significant). The qualitative picture of E1 thus carries over unchanged when jobs mix collectives, confirming that \Name's guarantees are collective-agnostic.

\begin{figure}[t]
\centering
\includegraphics[width=\columnwidth]{figure/fig7_e17_mixed.pdf}
\caption{Generality across collectives: P-attn vs.\ spine bandwidth under E17 (mixed AllReduce/AllGather/ReduceScatter/All-to-All, $n{=}10$). \Name leads at every scarce bandwidth; all policies converge above 800~Gbps.}
\label{fig:e17-mixed}
\end{figure}

\paragraph{Model-assumption robustness.}
The serial model~(\S\ref{sec:model}) is deliberately conservative. Varying overlap $\rho\in\{0,0.3,0.5,0.85,1.0\}$ leaves \Name stable ($82.5$--$87.5\%$ at 500\,Gbps, $92.5$--$97.5\%$ at 630\,Gbps) while other policies degrade; CASSINI, whose design assumes phase staggering, degrades most ($67.5\%{\to}45.0\%$ at 500\,Gbps) (Figure~\ref{fig:e11_overlap}).

\begin{figure}[t]
\centering
\includegraphics[width=\columnwidth]{figure/fig_e11_overlap.pdf}
\caption{Model-assumption robustness: P-attn vs.\ overlap factor $\rho$ across seven policies. (a)~500~Gbps spine; (b)~630~Gbps spine.}
\label{fig:e11_overlap}
\end{figure}

\subsection{Adversarial Workload}
\label{sec:eq2}

We next test whether \Name's advantage is specific to CRUX's favorable workload, or structural. In E2$'$ (inverted mapping) at 500~Gbps, \Name achieves $85.6\%$ P-attn against CRUX's $64.4\%$, a $21.1$~pp gap (Figure~\ref{fig:e2-killfield}). At 630~Gbps, \Name reaches $92.2\%$ while every baseline remains $\le 85.6\%$; at 800~Gbps \Name and DF tie at $96.7\%$. In E2-pro (favorable control), \Name is the only policy at $98.0\%$ at both 630 and 800~Gbps, with DF trailing at $90.0\%/98.0\%$ and CRUX at $86.0\%/90.0\%$. The gap under inverted mapping confirms \Name's advantage is structural and workload-agnostic.

\begin{figure}[h]
\centering
\includegraphics[width=\columnwidth]{figure/fig3_e2_orthogonal.pdf}
\caption{E2 orthogonal workload: (a)~E2$'$ (inverted), (b)~E2-pro (favorable).}
\label{fig:e2-killfield}
\end{figure}

\subsection{Dynamic Demand Reversal}
\label{sec:eq3}

We then evaluate whether \Name can maintain SLO guarantees when tier assignments change mid-run. In both control (E3, 800~Gbps) and kill (E3$'$, 630~Gbps) arms, \Name maintains $100.0\pm0.0\%$ P-attn across all windows with zero variance (Figure~\ref{fig:eval-e3}). In the control arm the best baselines in the final window are DF ($96.7\pm6.7\%$) and LL-S ($60.0\pm13.3\%$); in the kill arm DF collapses to $25.0\%$ by the final window, and CRUX's strict-priority isolation keeps standard jobs starved (S-cont $0.66$) even though it only reaches $47.5\%$ premium attention. \Name's stateless recomputation enables this: per-job targets are recomputed instantly from the new tier assignment, so the system transits without convergence delay.

\begin{figure}[h]
\centering
\includegraphics[width=\columnwidth]{figure/fig1_hero.pdf}
\caption{Dynamic tier swap: sliding-window P-attn in (a)~E3 control at 800~Gbps, (b)~E3$'$ kill at 630~Gbps. Gray bands W1--W3; dashed line marks swap.}
\label{fig:eval-e3}
\end{figure}

\textbf{Why the deficit-feedback controller fails.}\label{sec:eq4}
DF's kill-arm collapse is mechanistic: its scalar deficit $\pi$ cannot encode tier isolation and per-job capping simultaneously (\S\ref{sec:why-closed-form}). After the swap $\pi$ spans $-0.51$ to $+0.80$ across the 13 active jobs, scattering them across DSCP classes whose weights span $64\times$ ($\{1,2,4,8,16,32,64\}$), with a further $\exp(\pi{\cdot}K)$ share (clipped at $10\times$) inside each class---two properties \Name's deficit-based scheduling expresses directly and jointly. \Name exhibits zero variance across all seeds and windows (Figure~\ref{fig:d1-transient}).

\subsection{Trace-Driven Validation on the Lingjun 2023 Dataset}
\label{sec:trace}

We replay the densest 6-hour submission window of the Alibaba Lingjun 2023 dataset~\cite{LingjunTrace}, compressed onto the 600\,s horizon while preserving empirical inter-arrival distribution. This leaves 14 concurrent network jobs (9$\times$LLaMA-2-65B, 3$\times$GPT-3B, 2$\times$LLaMA-2-7B) on a 16-node FatTree at 400\,Gbps spine.

\Name achieves $75.0\pm11.2\%$ P-attn---statistically indistinguishable from the top baseline LL-S ($85.0\pm18.3\%$, $p{=}0.20$). The near-tie hides a systematic trade-off: LL-S buys its higher mean with per-seed variance ($18.3$~pp vs.\ $11.2$~pp), the worst worst-case SAS among the top two ($0.523{\pm}0.609$ vs.\ \Name's $0.731{\pm}0.269$), premium starvation in half of the seeds ($5/10$ vs.\ $1/10$; the mean-only plot in Figure~\ref{fig:trace-compare} obscures this tail cost), and $3\%$ fewer completed iterations ($2192.6$ vs.\ $2262.3$). Every other baseline trails by $26.4$~pp or more: SRPT $48.6\%$, DF $40.0\%$, Fair $35.7\%$, CRUX $34.3\%$, and CASSINI $20.7\%$.

\begin{figure}[h]
\centering
\includegraphics[width=\columnwidth]{figure/fig6_trace_compare.pdf}
\caption{Lingjun 2023 trace replay P-attn at 400\,Gbps ($n{=}10$). Asterisk marks $p<0.05$ vs.\ LongLiu.}
\label{fig:trace-compare}
\end{figure}

Table~\ref{tab:trace-replay} in Appendix~\ref{sec:app-supp} gives per-policy SAS breakdowns; the key finding is that \Name achieves both the highest worst-case SAS ($0.731\pm0.269$) and the highest aggregate throughput ($2262.3$ vs.\ $2192.6$ for LL-S), i.e., the near-tie in mean P-attn does not come at the cost of the tail.

\subsection{Hardware Validation}
\label{sec:eq5}

We validate on physical hardware in two ways: an NCCL AllReduce benchmark on a 50\,Gbps RoCEv2 testbed~(\S\ref{sec:impl-testbed}, primary), and a mechanism-boundary characterization.

\textbf{Physical testbed: NCCL AllReduce mechanism validation.}

Two-node dumbbell testbed (50\,Gbps RoCEv2, \S\ref{sec:impl-testbed}) with 6\,Gbps iperf3 UDP background. Each epoch issues fixed AllReduce operations on a 1024\,MB tensor. A staggered job pair swaps tightness at epoch~7; each job's baseline established via Phase-0 solo calibration. \Name{} (dynamic P6$\leftrightarrow$P2) vs.\ CRUX (static P3) across ten rounds. Over epochs 7--11, \Name{} achieves mean slowdown $1.062\times$ (95\%~CI $[1.024,1.099]$) vs.\ CRUX $1.520\times$ ($[1.476,1.564]$), a $30.1\%$ reduction ($t{=}20.31$, $p<0.0001$) across the six current-configuration rounds (Figure~\ref{fig:testbed-v6}). Across all ten rounds, the advantage is $22.6\%$ ($[14.9,30.4]$); excluding one Grubbs-confirmed outlier raises it to $25.1\%$ ($[19.0,31.2]$).

\begin{figure}[h]
\centering
\includegraphics[width=\columnwidth]{figure/fig6_testbed.pdf}
\caption{Physical testbed: (a)~epoch-by-epoch slowdown ($n{=}6$); (b)~per-round advantage with 95\%~CI.}
\label{fig:testbed-v6}
\end{figure}

\textbf{Mechanism boundary at the wire.}

DSCP-based strict-priority arbitration changes outcomes only when the wire is the bottleneck. A raw-RDMA probe confirms the switch realizes a $253\times$ priority service ratio, but with total demand ($18$--$37$\,Gbps) below the 50\,Gbps link, the scheduler helps the tight job modestly ($+3.5\%$ at 15\,Gbps, $+14.9\%$ at 30\,Gbps). This boundary is a deployment reality: \Name's guarantees hold exactly where switch egress is contended. Appendix~\ref{sec:app-boundary} gives the full trajectory and GPT-2 validation.

\begin{figure}[h]
\centering
\includegraphics[width=0.9\columnwidth]{figure/fig7_v6_summary.pdf}
\caption{Mean slowdown by background rate and SLO tightness. LongLiu gains $+3.5$\,pp at 15\,Gbps and $+14.9$\,pp at 30\,Gbps on tight jobs.}
\label{fig:v6-boundary-sum}
\end{figure}

\subsection{Summary}
\label{sec:eval-summary}

The evidence converges: \Name stratifies the field under scarcity~(\S\ref{sec:eq1}), its advantage is structural~(\S\ref{sec:eq2}), it sustains every premium SLO through immediate tier reversal~(\S\ref{sec:eq3}), its closed-form allocation avoids the expressiveness wall that collapses feedback controllers~(\S\ref{sec:eq4}), real traces confirm the results~(\S\ref{sec:trace}), and hardware validates the enforcement path while revealing the mechanism's boundary~(\S\ref{sec:eq5}).
